\documentclass[11pt, a4paper]{article}
\usepackage[utf8]{inputenc}
\usepackage[ngerman]{babel}
\usepackage[T1]{fontenc}
\usepackage{amsmath, amssymb, amsfonts}
\usepackage{geometry}
\geometry{a4paper, left=2.5cm, right=2.5cm, top=3cm, bottom=3cm}
\usepackage{hyperref}
\usepackage{titlesec}
\usepackage{enumitem}
\usepackage{setspace}
\usepackage{booktabs}
\setstretch{1.15}
\titleformat{\section}{\large\bfseries\sffamily}{\thesection}{1em}{}
\titleformat{\subsection}{\normalsize\bfseries\sffamily}{\thesubsection}{1em}{}
\newcommand{\mvec}[1]{\boldsymbol{#1}}
\newcommand{\clc}{\mathcal{C}\ell_{3}(\mathbb{C})}
\newcommand{\st}[1]{{\small\sffamily[#1]}}

\begin{document}
\begin{center}
  \sffamily
  {\Large\textbf{Clifford-BEM: Zusammenfassung der Vorarbeiten}}\\[0.3cm]
  {\large Ergebnisse zu AP 1 (Theorie), AP 2 (Diskretisierung) und AP 3 (Kompression)}\\[0.3cm]
  {\small Stand: Arbeitspapiere AP1 v0.27, AP2 v0.3, AP3 v0.5; C++-Kern v0.31}
\end{center}

\section*{Vorbemerkung}
Dieses Dokument fasst die Ergebnisse der drei Arbeitspapiere zusammen. Die Statusmarken
bedeuten:
\begin{itemize}[leftmargin=*]
  \item \st{bewiesen}: vollständiger Beweis im Arbeitspapier.
  \item \st{exakt}: symbolische Identität, maschinell auf null gekürzt.
  \item \st{numerisch}: durch Rechnungen belegt.
  \item \st{Vermutung}: offen.
\end{itemize}
Alle Beweise sind vorläufig und vor einer Veröffentlichung unabhängig zu prüfen. Die
Literaturangaben in den Arbeitspapieren sind ungeprüft. Die numerischen Ergebnisse stammen
aus Python-Prototypen mit 3\,GB Arbeitsspeicher; das begrenzt die erreichbaren Netzgrößen.

\paragraph{Die wichtigsten Ergebnisse.}
\begin{enumerate}[leftmargin=*]
  \item Die Formulierung $T_1=E_2^++E_1^-J$ ist für alle passiven Medien frei von
    künstlichen Resonanzen, auch für verlustfreie Metalle. Die naheliegende Alternative
    $E_2^++J^{-1}E_1^-J$ ist es nicht.
  \item Mit der \glqq Standardwahl\grqq{} der Hilfsparameter ist das Problem auf dem vollen
    Multivektorraum eindeutig lösbar, für alle passiven Medien mit
    $\operatorname{Im}(\varepsilon\mu)\ge0$, chiral unter $\operatorname{Im}k_\pm^2\ge0$.
  \item An Kanten verliert jede $L^2$-Formulierung die Stabilität für Kontraste innerhalb
    zweier explizit bekannter komplexer Schleifen. Das Symbol von $T_1$ faktorisiert dort
    exakt in physikalische und triviale Hilfsblöcke. Gewichtete Räume stellen die
    Stabilität her, außer im Resonanzfenster.
  \item Am Würfel bestätigt der 3D-Prototyp die Kantenvorhersage quantitativ: beobachtete
    Rate $0{,}31$, vorhergesagt $0{,}32$. Mit Kanten- und Eck-Blockvorkonditionierung
    sind die Iterationszahlen von der Kantenauflösung unabhängig, auch mit
    $\mathcal{H}$-Matrix-Kompression.
  \item Auf der Kugel stimmt die Extinktion nach Extrapolation mit Mie überein: für Glas
    auf $4\cdot10^{-5}$ (bis 23\,040 Unbekannte), für Gold auf $0{,}4\,\%$.
\end{enumerate}

%=====================================================================
\section{AP 1: Theorie}
%=====================================================================

\subsection{Grundlagen}
\begin{itemize}[leftmargin=*]
  \item Die Maxwell-Gleichungen sind für Felder in den Graden 1 und 2 äquivalent zu
    $(\nabla-ik)\mvec{F}=0$ mit $\mvec{F}=\sqrt\varepsilon\mvec{E}+I\sqrt\mu\mvec{H}$.
    Beidseitige Fundamentallösung, Cauchy-Darstellung innen und außen. Die
    Strahlungsbedingung ergibt sich aus dem Beweis als Projektion
    $(1-\hat{\mvec{x}})\mvec{F}=o(1/r)$; es gilt $\mathcal{E}_\Gamma^2=1$. \st{bewiesen}
  \item Kontrollen: Cauchy-Formel auf der Kugel auf $10^{-15}$; Transmissionsabbildung $J$
    an exakten Fresnel-Lösungen auf $10^{-16}$, auch für Gold. \st{numerisch}
\end{itemize}

\subsection{Formulierung und Stabilität auf glatten Rändern}
\begin{itemize}[leftmargin=*]
  \item \emph{Hauptsymbol:} Fredholm vom Index 0 genau dann, wenn
    $\varepsilon_1\neq-\varepsilon_2$ und $\mu_1\neq-\mu_2$. Die Ausnahme ist das
    quasistatische Oberflächenplasmon. Der Operator ist elliptisch nullter Ordnung, aber
    für Kontrast $\neq1$ \emph{nicht} Identität plus kompakt. \st{bewiesen}
  \item \emph{Falsche Resonanzen der naheliegenden Gleichung:} Die Variante mit
    $W=J^{-1}$ hat Kerne bei Hohlraumplasmonen der komplementären Geometrie (Kugel:
    $\varepsilon_1^*\approx-0{,}6935$, Stabilitätswert $\to6\cdot10^{-5}$). \st{numerisch}
  \item \emph{Resonanzfreiheit von $T_1$:} Eine Flussidentität für Dirac-Felder macht das
    duale Problem für alle passiven Medien eindeutig lösbar. $T_1$ ist daher genau dann
    invertierbar, wenn das physikalische Problem eindeutig lösbar ist. Für
    verlustbehaftete Medien liegt das Symbolspektrum in der rechten Halbebene (Gold:
    $0{,}18+3{,}32i$ und $0{,}10+1{,}51i$). \st{bewiesen}
  \item \emph{Niederfrequenz:} kein Zusammenbruch auf Kugel und Torus bis $\omega a=10^{-6}$.
    Auf dem Torus springt der Modus $m=0$ bei exakt $\omega=0$; das ist der topologische
    Sektor der harmonischen Felder. \st{numerisch}
\end{itemize}

\subsection{Eindeutigkeit im vollen Multivektorraum}
\begin{itemize}[leftmargin=*]
  \item \emph{Standardwahl} ($a_{\mathrm H}=\sqrt{r_\mu}$, $b_{\mathrm H}=\sqrt{r_\varepsilon}$):
    Die Hilfspotentiale $\alpha/\sqrt\mu$, $\beta/\sqrt\varepsilon$ werden stetig. Die
    Maxwell-Struktur erzwingt stetige Normalableitungen, und die Hilfsanteile verschwinden
    per Rellich. Die Eindeutigkeit gilt für $\operatorname{Im}(\varepsilon_1\mu_1)\ge0$, also für
    alle nichtmagnetischen passiven Medien. \st{bewiesen, reguläre Lösungen}
  \item \emph{Chirale Innenmedien:} Dasselbe helizitätsweise, mit einer über die
    Helizitäten gewichteten Randform, unter $\operatorname{Im}k_\pm^2\ge0$. Das schließt
    Gold mit $|\chi|\lesssim0{,}18$ ein, nicht aber stark chirale Metalle. \st{bewiesen}
  \item Die alternative \glqq Energiewahl\grqq{} hat ebenfalls einen Eindeutigkeitsbeweis,
    ist aber an Kanten mit den nötigen gewichteten Räumen unverträglich (siehe 1.4).
    Empfohlen wird daher die Standardwahl.
  \item Numerische Bestätigung: Bei Streuung im vollen Raum liegen die Hilfsanteile auf
    Residuumsniveau, und die Extinktion stimmt auf 5 Stellen mit Mie überein.
\end{itemize}

\subsection{Kanten}
\begin{itemize}[leftmargin=*]
  \item \emph{Kritische Menge:} Das Konormalsymbol an einer Kante mit Winkel $\alpha$ ist
    genau dann singulär, wenn $(1+\kappa)/(1-\kappa)=\pm S(a)$ gilt, mit
    $S(a)=\sin((\pi-\alpha)a)/\sin(\pi a)$ und $\operatorname{Re}a=c$ ($c=\tfrac12$: $L^2$). Für
    $\alpha=\pi/2$ sind das in $L^2$ zwei komplexe Schleifen zwischen $-5{,}83$ und
    $-0{,}17$, die bis $\operatorname{Im}\kappa\approx2{,}1$ reichen. Die Endpunkte sind
    $-\cot^2(\alpha/4)$ und $-\tan^2(\alpha/4)$; im Energie-Grenzfall ergibt sich das
    bekannte Intervall $[-3,-\tfrac13]$. Innerhalb der Schleifen ist der Index $\neq0$.
    \st{exakt}
  \item \emph{Blockfaktorisierung:}
    $\det\sigma_M(T_1)=\tfrac1{256}\,\hat g(r_\varepsilon;a)\,\hat g(r_\mu;a)\,
    \hat g\big(\tfrac{a_{\mathrm H}}{\sqrt{r_\mu}};a-1\big)\,\hat g\big(\tfrac{b_{\mathrm H}}{\sqrt{r_\varepsilon}};a-1\big)$
    mit $\hat g(c;a)=(1-c)^2S(a)^2-(1+c)^2$,
    mit einem TM-, einem TE- und zwei Hilfsblöcken, exakt für allgemeines $\alpha$. Bei
    Standardwahl sind die Hilfsblöcke für jedes Gewicht trivial. Bei Energiewahl trägt ein
    Hilfsblock die Windung $-1$; in $L^2$ kompensiert das die physikalische Windung $+1$,
    im gewichteten Raum ist die Gesamtwindung dagegen $-1$. \st{exakt}
  \item \emph{Gewichte:} $L^2(r^{2\beta})$ mit $\beta>\tfrac12-\min_j\operatorname{Re}a_j$, mit den
    Singulärexponenten $a_j$. Einordnung für $\alpha=\pi/2$:
\end{itemize}
\begin{center}\small
\begin{tabular}{lccc}
\toprule
$\kappa$ & in $L^2$-Schleife & $\min\operatorname{Re}a_j$ & nötiges $\beta$\\
\midrule
Gold/Vakuum $-11+1{,}2i$ & nein & 0,59 & 0\\
Silber/Vakuum $-15+0{,}5i$ & nein & -- & 0\\
Gold/Wasser $-6{,}2+0{,}68i$ & knapp nein & 0,52 & 0\\
$-4$ (verlustfrei) & ja & 0,37 & $\ge0{,}13$\\
$-3+0{,}3i$ & ja & 0,19 & $\ge0{,}32$\\
$\kappa\in[-3,-\tfrac13]$ verlustfrei & ja & 0 & keines\\
\bottomrule
\end{tabular}
\end{center}

\subsection{Kegelspitzen}
\begin{itemize}[leftmargin=*]
  \item Physikalische kritische Kontraste
    $\kappa(\nu,m)=-P_\nu^{m\prime}(-c)P_\nu^m(c)/(P_\nu^{m\prime}(c)P_\nu^m(-c))$ mit
    $c=\cos\theta_0$; der Randwert in $L^2$ ist $-\cot^2(\theta_0/2)$. \st{exakt}
  \item Das Konormalsymbol von $T_1$ reproduziert diese Kontraste in 54 Fällen auf
    $9\cdot10^{-13}$; die Hilfsgrade erzeugen keine weiteren. \st{numerisch}
  \item An spitzen Kegeln ($\theta_0=20^\circ$) reicht die $L^2$-Menge bis $-32$. Gold und
    Silber in Vakuum liegen dann innerhalb und brauchen $\beta\approx0{,}39$ bzw. $0{,}47$,
    nahe der Muckenhoupt-Grenze $\tfrac12$.
\end{itemize}

\subsection{Chirale Grenzflächen}
\begin{itemize}[leftmargin=*]
  \item $J$ mischt über die Materialmatrix $C_j$ den vektoriellen und den bivektoriellen
    Normalanteil; an Fresnel-Lösungen auf $6\cdot10^{-16}$ verifiziert. \st{numerisch}
  \item \emph{Glatter Rand:} kritisch ist $(\varepsilon_1+1)(\mu_1+1)=\chi_1^2$ (Vakuum außen);
    mit chiralem Außenmedium $(\varepsilon_1+\varepsilon_2)(\mu_1+\mu_2)=(\chi_1+\chi_2)^2$. \st{exakt / numerisch}
  \item \emph{Kanten:}
    $\det=\frac{\varepsilon\mu}{16(\varepsilon\mu-\chi^2)^2}(\mathcal A-\mathcal BS^2+\mathcal CS^4)$
    mit ganzzahligen Polynomkoeffizienten, auf 40 Stellen bestätigt. Die Verschiebung ist
    von der Ordnung $\chi^2$; für starke Chiralität reicht die kritische Menge bis zu
    positiven $\varepsilon$. \st{numerisch}
\end{itemize}

\subsection{2D-Prototyp (Quadrat)}
\begin{itemize}[leftmargin=*]
  \item In $L^2$ fällt der kleinste Singulärwert innerhalb der Schleifen mit der
    vorhergesagten Rate (gemessen $0{,}111$, vorhergesagt $0{,}127$). Auf geometrischen Netzen
    bleibt er in $L^2(r^{2\beta})$ beschränkt. \st{numerisch}
  \item Eck-Blockvorkonditionierung: etwa 14 GMRES-Iterationen, unabhängig von der
    Eckauflösung, statt bis zu 60.
  \item \emph{Resonanzfenster:} Resonanzen in $[-3,-\tfrac13]$ mit
    $\operatorname{Im}\kappa\lesssim1$ konvergieren mit Kollokation nicht: Eckwellen $r^{-i\tau}$ (in $\ln r$
    etwa 11 Dekaden Wellenlänge) werden am innersten Element reflektiert.
\end{itemize}

\subsection{Resonanzfenster: Lösung in 2D, Befunde in 3D}
\begin{itemize}[leftmargin=*]
  \item \emph{2D gelöst:} Galerkin mit komplexer Streckung $r\mapsto r\,(r/r_0)^{i\theta}$ ($\theta>0$,
    $c+\tau\theta>\tfrac12$) und Konturtestung. Tiefenunabhängig auf $10^{-7}$, konvergent, $\theta$-unabhängig
    auf $6\cdot10^{-4}$, auch verlustfrei. Mit falschem Vorzeichen von $\theta$ konvergiert das Verfahren gegen eine
    unphysikalische Lösung. \st{numerisch}
  \item \emph{3D-Kanten:} Die Streckung erzeugt Nullstellen von $z\cdot z$ auf dem komplexen Rand (Koordinate
    entlang der Kante bleibt reell); begrenzte Phase dämpft zu wenig. Angereicherte Eckelemente (auch mit der
    exakten Mellin-Mode) und ein Absorber über elementweises $\operatorname{Im}\kappa$ sind nicht transparent.
    An Spitzen ist eine isotrope Streckung zulässig. \st{numerisch}
  \item \emph{Abgerundete Kanten:} Mit Rundungsradius $\rho$ konvergiert die 3D-Rechnung im Fenster monoton, hängt
    aber stark von $\rho$ ab (etwa Halbierung der Extinktion von $\rho=0{,}4$ auf $0{,}2$). \st{numerisch}
\end{itemize}

%=====================================================================
\section{AP 2: 3D-Galerkin-Diskretisierung}
%=====================================================================
\begin{itemize}[leftmargin=*]
  \item \emph{Verfahren:} ebene Dreiecke, stückweise konstante $\clc$-Dichten, orthonormale
    Basis. Das Nahfeld wird analytisch integriert; der Selbstterm des Cauchy-Anteils ist
    wegen Antisymmetrie exakt null.
  \item \emph{Spurtrennung:} Konvergenz erster Ordnung (Fehler $4{,}6\,\%\to1{,}4\,\%$ von 80 auf
    320 Dreiecke), für reelles und komplexes $k$.
  \item \emph{Kugel gegen Mie} (dicht, bis 500 Dreiecke, extrapoliert): Glas $0{,}24\,\%$, Gold
    $0{,}47\,\%$, nahe Plasmonresonanz $0{,}73\,\%$. Der Hilfsanteil sinkt mit $h$.
  \item \emph{Stabilität auf der Kugel:} $\sigma_{\min}(T_1)$ ist beschränkt (Glas $\approx0{,}45$,
    Gold $\approx0{,}21$). GMRES: Glas 12--14, Gold 29--37 mit sättigendem Anstieg.
  \item \emph{Würfel} (Netz geometrisch zu den Kanten gradiert, Zerlegung in 8
    Symmetriesektoren, bis 13\,824 Unbekannte):
\end{itemize}
\begin{center}\small
\begin{tabular}{lccc}
\toprule
$\kappa$ & beobachtete $L^2$-Rate von $\sigma_{\min}$ & Vorhersage Kante / Vergleichskegel & Bewertung\\
\midrule
Glas & $0{,}001$ & -- & stabil\\
Gold & $0{,}06$, fallend & -- / -- & stabil (Tendenz)\\
$-3+0{,}3i$ & $0{,}31$ & $0{,}32$ / $0{,}34$ & Kante bestätigt\\
$-4+0{,}3i$ & $0{,}28$ & $0{,}12$ / $0{,}39$ & Ecken dominieren\\
\bottomrule
\end{tabular}
\end{center}
\begin{itemize}[leftmargin=*]
  \item \emph{Gewichtete Norm am Würfel:} Das Kantengewicht bremst den Abfall deutlich. An
    den Ecken reicht es nicht; dort ist ein eigenes Eckgewicht nötig.
  \item \emph{Kanten- und Eck-Blockvorkonditionierung:} Die Iterationen sind ab $L=3$
    konstant, für alle Kontraste:
\end{itemize}
\begin{center}\small
\begin{tabular}{lccc}
\toprule
$\kappa$ & ohne ($L=1\to5$) & $R=0{,}25$ & $R=0{,}5$\\
\midrule
Glas & $11\to12$ & 10 & 9\\
Gold & $43\to56$ & 33 & 27\\
$-4+0{,}3i$ & $42\to67$ & 36 & 27\\
$-3+0{,}3i$ & $41\to72$ & 36 & 28\\
\bottomrule
\end{tabular}
\end{center}
Nötig ist ein fester physikalischer Radius $R\gtrsim\tfrac18$; punktweise Diagonalblöcke
allein helfen nicht.

%=====================================================================
\section{AP 3: Kompression}
%=====================================================================
\begin{itemize}[leftmargin=*]
  \item \emph{Ansatz:} Gespeichert werden die vier Kernkomponenten je Dreieckspaar statt der
    $8\times8$-Blöcke, was schon Faktor 16 spart. Dazu kommen ein Clusterbaum, ACA mit
    partieller Pivotisierung und QR/SVD-Nachkompression. Der Fehler des
    Matrix-Vektor-Produkts ist $\approx0{,}1\,\varepsilon$.
  \item \emph{Gemeinsame gegen komponentenweise ACA} (H3, erster Teil): 10--19\,\% weniger
    Speicher bei gleicher Genauigkeit, dazu schnellerer Aufbau und schnelleres
    Matrix-Vektor-Produkt.
  \item \emph{Skalierung} (H3, zweiter Teil): Im Python-Prototyp (bis 3920 Dreiecke) noch
    präasymptotisch ($N^{1{,}39}$). Mit dem C++-Kern bis 20\,480 Dreiecke (163\,840 Unbekannte):
    Speicher $\mathcal{O}(N\log^2N)$ für Kugel und Würfel mit Kanten, der Quotient ist auf
    $\pm7\,\%$ konstant. Die Frequenz kostet bei fester Auflösung je Wellenlänge bis $kD\approx16$ nur
    $10\,\%$. Stark gradierte Netze sind wegen der Nahfeldkriterien derzeit dicht-dominiert.
  \item \emph{Kugel jenseits der dichten Grenze} (bis 23\,040 Unbekannte):
\end{itemize}
\begin{center}\small
\begin{tabular}{lcccc}
\toprule
& GMRES & $Q_{\mathrm{ext}}$ bei 2880 Dreiecken & extrapoliert & Mie\\
\midrule
Glas, $\omega a=1$ & 10 (konstant) & 0,21371 & 0,21511 & 0,21510\\
Gold, $\omega a=0{,}5$ & 25 (konstant) & 0,60764 & 0,59247 & 0,59002\\
\bottomrule
\end{tabular}
\end{center}
\begin{itemize}[leftmargin=*]
  \item \emph{C++-Streulöser} (v0.2): $T_1$ mit $\mathcal{H}$-Operatoren und GMRES bis 92\,160 Unbekannte.
    Mit Sauter-Schwab-Quadratur (v0.4) konvergieren Glas und Gold mit Ordnung 2; nach Extrapolation
    liegen beide auf $10^{-5}$ bis $10^{-4}$ bei Mie. Die langsamere Konvergenz bei Gold lag an der
    Nahfeldquadratur. Die Iterationen sind von $N$ unabhängig (10 bzw.\ 24).
  \item \emph{Kombinierter Löser am Würfel:} anisotropes, nichtkonformes Kantennetz,
    $\mathcal{H}$-Operatoren und Kanten-/Eckblöcke. Die Iterationen bleiben konstant (49 für
    $-3+0{,}3i$, 47 für Gold, $L=2\dots5$), ohne Blöcke wachsen sie ($73\to106$ bzw.\ $64\to71$).
    Die Extinktion fällt für Gold monoton und schwankt für $-3+0{,}3i$, passend zum
    Resonanzfenster.
  \item \emph{Chirale Medien} (v0.6): $Q_\pm$ und CD treffen die chirale Mie-Lösung auf $10^{-4}$.
  \item \emph{Mehrere Körper, Gmsh, Spektren} (v0.7, v0.8): Born-Kuhn-Dimer mit bisignatem, orientierungsgemitteltem
    CD-Spektrum; Enantiomere mit entgegengesetztem CD auf $10^{-5}$; Goldkugel in Wasser gegen Mie.
  \item \emph{Vorkonditionierung} (v0.9, v0.10): Kanten-/Eckblöcke auch für mehrere Körper und chirale Medien;
    hierarchische Faktorisierung (HODLR) als direkter Löser bzw.\ Vorkonditionierer, lohnend bei vielen rechten
    Seiten, an Plasmonresonanzen nur begrenzt.
  \item \emph{Nahfeld auf gestreckten Elementen} (v0.5): Sauter-Schwab ist dort ungenau (bis 1,4\,\%); halbanalytische,
    gradierte Regel.
  \item \emph{Silberwürfel} (50\,nm, Wasser): Die Hauptresonanz verschiebt sich mit schärferen Kanten nach Rot
    (grob 440, 460, 470\,nm für 10, 7,5, 5\,nm Rundungsradius).
  \item \emph{Beschichtete Grenzflächen} (v0.13): exakte Schichtformulierung für verschachtelte Gebiete (Kern-Schale,
    Mehrfachschichten, chirale Schichten); Konsistenz und Fredholm-Eigenschaft bewiesen, Eindeutigkeit vermutet.
    Beschichtete Kugel gegen Aden-Kerker mit Ordnung~2 in Extinktion und Phase der Vorwärtsamplitude. Bei Schichten
    dünner als die Elemente trägt die Rechnung einen systematischen Fehler; die Schichtwirkung als Differenz zu einer
    neutralen Rechnung (gleiche Netze, Schicht aus Außenmedium) ist bis $d/h\approx0{,}04$ auf wenige Prozent genau.
    Silberkugel (20\,nm, Wasser) mit 2\,nm Oxid ($n=1{,}7$): Rotverschiebung der Dipolresonanz um etwa 10\,nm,
    Phasenänderung der Vorwärtsamplitude bis 0,9\,rad, beides gegen Aden-Kerker auf 0,5--7\,\% an der Resonanz.
    Abgerundeter Silberwürfel (50\,nm) mit derselben Schicht: Rotverschiebung um etwa 20\,nm, Phasenänderung bis
    0,5\,rad.
    Dünnschicht-Näherung erster Ordnung auf einer Fläche (v0.14): für $d\ll h$ genauer als die Zwei-Flächen-Rechnung,
    besonders in der Phase ($d/a\lesssim0{,}03$), bei etwa einem Fünftel der Kosten; Modellfehler $\approx1{,}5\,d/a$.
    In zweiter Ordnung (v0.15, Dirac-Form mit Krümmung) trifft sie die Schichtwirkung bis $d/a=0{,}05$ auf 0,1--0,7\,\%
    und die Silberkugel mit 2\,nm Oxid so genau wie die exakte Rechnung. Seit v0.16 mit punktweise konsistenten zweiten
    Ableitungen und für mehrere Körper; sie braucht eine auf dem Netz aufgelöste Krümmung. Seit v0.17 auch für chirale
    Schichten: Der Zirkulardichroismus einer dünnen chiralen Hülle konvergiert mit Ordnung~2 gegen eine neue Mie-Lösung
    für geschichtete chirale Kugeln. Anwendung (v0.18): Goldkugel mit 1\,nm chiraler Molekülhülle in Wasser;
    der plasmoninduzierte CD wird an der Resonanz auf 2--7\,\% getroffen. Die Schicht als Zweitor (v0.19, nach dem
    Vorbild der S-Matrix) ist stabil auch für Schichten dicker als die Netzelemente und bis $d/a\approx0{,}1$ genauer als
    die Diskretisierung ohne Schicht. Seit v0.20 mit Mehrfachschichten und automatischer Unterteilung dicker Schichten
    konvergent bis $d/a=0{,}5$; der CD einer chiralen Schicht im Abstand 0--8\,nm zu einer Goldkugel wird bis auf einen
    konstanten Diskretisierungsversatz getroffen. Seit v0.21 auch chirale Außenmedien (Teilchen in chiraler Lösung),
    validiert gegen eine Mie-Lösung mit chiralem Außengebiet. Seit v0.22 rechnet das Zweitor mehrere Körper, etwa Dimere mit
    chiraler Molekülschicht im Spalt. Seit v0.24 Nahfeldkarten: Im Spalt des Dimers ist die optische Chiralität bis zehnfach
    verstärkt. Seit v0.26 Dipolanregung: Zerfallsraten und Quantenausbeute von Fluorophoren vor Nanostrukturen, seit v0.27 die
    Fluoreszenzverstärkung (im Spalt eines Gold-Dimers etwa 260-fach).
    \st{numerisch}
\end{itemize}

%=====================================================================
\section{Stand der Hypothesen}
%=====================================================================
\begin{description}[font=\normalfont\bfseries, style=nextline, leftmargin=1.2em]
  \item[H1 (Stabilität):] Auf glatten Rändern gestützt: beschränktes $\sigma_{\min}$,
    Iterationen konstant bzw.\ schwach und sättigend wachsend. Die punktweise
    Vorkonditionierung $2(1+J)^{-1}$ bringt dabei wenig. An Kanten und Ecken gestützt in der
    revidierten Form mit Kanten- und Eck-Blockvorkonditionierung. Ausdrücklich ausgenommen
    ist das Resonanzfenster.
  \item[H2 (Diskretisierung):] Auf der Kugel stabil. Am Würfel verhält sich die Diskretisierung
    wie vorhergesagt: instabil in $L^2$ innerhalb der Schleifen, in der gewichteten Norm
    teilweise stabilisiert, wobei ein Eckgewicht fehlt. Ein Stabilitätsbeweis
    (Gårding-Ungleichung) fehlt.
  \item[H3 (Kompression):] Der erste Teil ist in der Form „Multivektor-Pivots senken den
    Speicher“ widerlegt: Eine Kreuzapproximation über der Algebra braucht bei gleicher Genauigkeit
    2,3- bis 2,9-mal mehr Speicher als die gemeinsame skalare ACA, weil der Kern skalar separabel ist.
    Gestützt ist die gemeinsame Behandlung der Grade mit skalaren Zeilen- und Paravektor-Spaltenfaktoren
    (etwa 19\,\% weniger Niedrigrang-Speicher als komponentenweise). Der zweite Teil ist bis etwa 160\,000 Unbekannte und $kD\lesssim16$
    gestützt: $\mathcal{O}(N\log^2N)$. Offen sind stark gradierte Netze und elektrisch große
    Streuer.
\end{description}

%=====================================================================
\section{Offene Punkte}
%=====================================================================
\begin{itemize}[leftmargin=*]
  \item \emph{Theorie:} Fredholm-Theorie auf Lipschitz-Rändern (Kantenoperatoren entlang der
    Kante, Spurregularität); echte Würfelecke (Mellin-Symbol auf dem sphärischen Dreieck);
    Beweis der Übereinstimmung an Kegelspitzen; stark chirale Metalle; Tellegen-Medien.
  \item \emph{Resonanzfenster:} transparenter Abschluss an 3D-Kanten (asymptotische diskrete DtN aus der
    Schichtstruktur, Hardy-Raum-Ansätze); Streckung um Spitzen im C++-Kern; unabhängige Referenz im Fenster.
  \item \emph{Diskretisierung:} Gårding-Ungleichung für $T_1$; ein kombiniertes Gewicht für Kanten und Ecken;
    Nahfeldkriterien für gestreckte Elemente; Sauter-Schwab für nichtkonforme Netze.
  \item \emph{Kompression und Löser:} ACA+; H-LU mit starker Zulässigkeit; Krylov-Recycling für viele rechte Seiten;
    Parallelisierung und Tests auf Mehrkernrechnern.
  \item \emph{Modellierung:} Substrate, gekrümmte Elemente, Ansätze höherer Ordnung; Eindeutigkeit und
    Vorkonditionierung für geschichtete Körper.
  \item \emph{Formales:} unabhängige Prüfung aller Beweise; Prüfung der Literatur.
\end{itemize}

\end{document}
